\documentclass[a4paper,11pt,oneside]{book}
% preamble_build.tex: copy of Maximus's preamble.tex for the DE v3 build.
% The blocks marked [BUILD COPY CHANGE] are the earlier build fixes.
% The blocks marked [V3 ADDITION] are new in this version (example box,
% spacing). Everything else is unchanged from preamble.tex.

% For margins
\usepackage[a4paper, margin = 20mm]{geometry}

\usepackage[Lenny]{fncychap}
% Rejne, Glenn, Lenny

\usepackage{etoc}

% [V3 ADDITION] cancel: used to show the dx's "cancelling" in Section 11.3
\usepackage{cancel}
\renewcommand{\CancelColor}{\color{red}}

% For enumeration & math symbols
\usepackage{enumerate}
\usepackage{amsmath,amsfonts,amssymb,amsthm,graphicx,lipsum}

% For diagrams, plots and boxes
\usepackage{pgfplots}
\pgfplotsset{compat=1.18}
\usepackage{tikz}
\usepackage[most]{tcolorbox}
\tcbuselibrary{skins}
\usepackage{xcolor}
\pgfplotsset{
    standard/.style={clip = false,
    axis line style = thick,
    trig format=rad,
    axis lines = middle,
    every axis x label/.style={at={(current axis.right of origin)},anchor= west},
    every axis y label/.style={at={(current axis.above origin)},anchor=south}
    }
}

% For GC font
\usepackage{fontspec}
% [BUILD COPY CHANGE 1] The TI-84 key font is loaded only if the file exists
% (checked in ./Fonts/... and ../Fonts/...). Otherwise \ti{} falls back to a
% small boxed sans-serif key. Known letters of the TI font are mapped to their
% key names in the fallback: o = [y=], z = [mode], f = [L3].
\newcommand{\tikey}[1]{\tcbox[on line, enhanced, boxsep=0pt, left=2pt, right=2pt,
  top=1pt, bottom=1pt, boxrule=0.5pt, arc=1.5pt, colback=white,
  colframe=black!70]{\sffamily\footnotesize\strut #1}}
\ExplSyntaxOn
\NewDocumentCommand{\tifallback}{m}{
  \str_case:nnF {#1}
    { {o}{\tikey{y=}} {z}{\tikey{mode}} {f}{\tikey{L3}} }
    { \tikey{#1} }
}
\ExplSyntaxOff
\def\tifontcase{0}
\IfFileExists{Fonts/TI-84 Plus CE Key Fonts/TI84PlusCEKeys.ttf}
  {\def\tifontcase{1}}{}
\IfFileExists{../Fonts/TI-84 Plus CE Key Fonts/TI84PlusCEKeys.ttf}
  {\def\tifontcase{2}}{}
\ifcase\tifontcase\relax
  % 0: font not found, use boxed fallback keys
  \newcommand{\ti}[1]{\tifallback{#1}}
\or
  % 1: original location (as in preamble.tex)
  \newfontfamily\tifont[
    Path = {Fonts/TI-84 Plus CE Key Fonts/},
    UprightFont = {TI84PlusCEKeys.ttf}
  ]{}
  \newcommand{\ti}[1]{{\tifont #1}}
\or
  % 2: one level up (when compiling inside draft-ch1/)
  \newfontfamily\tifont[
    Path = {../Fonts/TI-84 Plus CE Key Fonts/},
    UprightFont = {TI84PlusCEKeys.ttf}
  ]{}
  \newcommand{\ti}[1]{{\tifont #1}}
\fi

% [FONT SETUP copied from h2-math-notes/aoi/notes/preamble_build.tex (DI notes)]
% [AoI items 17/18/37] GC DISPLAY FONT. Every GC key, command, menu item and
% screen line goes through \gcfont, so the font can be swapped HERE in one place.
%   \gc{fnInt(}      command / menu item / stored variable in the display font
%   \gckey{2nd}      a calculator key (display font in a small key outline)
%   \gcscreen{...}   a calculator screen (display font on a grey-green LCD panel)
% The font file lives in ../Fonts (shared by notes, slides and worksheet).
\def\gcfontfile{VT323-Regular.ttf}
\IfFileExists{Fonts/\gcfontfile}{%
  \newfontfamily\gcfont[Path=Fonts/, Scale=1.32]{\gcfontfile}}{%
  \newcommand\gcfont{\ttfamily}}
\newcommand{\gc}[1]{\mbox{\gcfont #1}}
\newcommand{\gckey}[1]{\tcbox[on line, enhanced, boxsep=0pt, left=2pt, right=2pt,
  top=0.6pt, bottom=0.2pt, boxrule=0.5pt, arc=1.5pt, colback=black!4,
  colframe=black!65]{\gcfont\strut #1}}
\definecolor{lcd}{RGB}{214,222,204}
\newcommand{\gcscreen}[1]{\tcbox[on line, enhanced, boxsep=1pt, left=4pt, right=4pt,
  top=2pt, bottom=2pt, boxrule=1.2pt, arc=1pt, colback=lcd, colframe=black!70]{%
  \gcfont\begin{tabular}{@{}l@{}}#1\end{tabular}}}
% integral template as shown on the TI-84 Plus CE (MathPrint): \gcint{0}{2}{e^(-X²)}{X}
\newcommand{\gcint}[4]{{\gcfont\raisebox{-0.35ex}{\Large ∫}\hspace{-0.15em}%
  \raisebox{-0.2ex}{\begin{tabular}{@{}l@{}}\footnotesize #2\\[-1.3ex]\footnotesize #1\end{tabular}}%
  \hspace{0.1em}(#3)d#4}}
% old key macro kept as an alias so nothing breaks
\renewcommand{\ti}[1]{\gckey{#1}}

% SI units
\usepackage{siunitx}

% setspace
%\usepackage{setspace}
%\onehalfspacing

% For hyperlinks within and out of document
\usepackage{hyperref}
\hypersetup{
    colorlinks = true,
    linkcolor = black,
    urlcolor = blue,
}   

% For captions for pics
\usepackage{capt-of}


% tcolorbox custom environments
% Definition ========================================================
\newtcolorbox[auto counter,number within=section]{definition}[2][drop shadow]{%
enhanced, before skip=5mm, after skip=6mm, % [V3 ADDITION] spacing
colback=blue!5!white,colframe=blue!50!black,fonttitle=\bfseries,
title= Definition~\thetcbcounter. #2,#1}

% Theorem ==========================================================
\newtcolorbox[auto counter,number within=section]{theorem}[2][drop shadow]{%
enhanced, before skip=5mm, after skip=6mm, % [V3 ADDITION] spacing
colback=black!5!white,colframe=black,fonttitle=\bfseries,
title= Theorem~\thetcbcounter. #2,#1}

% Caution box ======================================================
\newtcolorbox{caution}[1][]{enhanced,
before skip=5mm,after skip=6mm, % [V3 ADDITION] was 2mm/3mm
boxrule=1pt,left=11mm,right=2mm,top=3mm,bottom=3mm,
colback=yellow!50,
colframe=yellow!20!black,
sharp corners,rounded corners=southeast,arc is angular,arc=3mm,
underlay={%
\path[fill=tcbcolback!80!black] ([yshift=3mm]interior.south east)--++(-0.4,-0.1)--++(0.1,-0.2);
\path[draw=tcbcolframe,shorten <=-0.05mm,shorten >=-0.05mm] ([yshift=3mm]interior.south east)--++(-0.4,-0.1)--++(0.1,-0.2);
\path[fill=yellow!50!black,draw=none] (interior.south west) rectangle node[white]{\Huge\bfseries !} ([xshift=9mm]interior.north west);
},
drop fuzzy shadow,#1}

% Question box ======================================================
\newtcolorbox{question}[1][]{enhanced,
before skip=5mm,after skip=6mm, % [V3 ADDITION] was 2mm/3mm
boxrule=1pt,left=11mm,right=2mm,top=3mm,bottom=3mm,
colback=blue!20!white,
colframe=blue!20!black,
sharp corners,rounded corners=southeast,arc is angular,arc=3mm,
underlay={%
\path[fill=tcbcolback!80!black] ([yshift=3mm]interior.south east)--++(-0.4,-0.1)--++(0.1,-0.2);
\path[draw=tcbcolframe,shorten <=-0.05mm,shorten >=-0.05mm] ([yshift=3mm]interior.south east)--++(-0.4,-0.1)--++(0.1,-0.2);
\path[fill=blue!50!black,draw=none] (interior.south west) rectangle node[white]{\Huge\bfseries ?} ([xshift=9mm]interior.north west);
},
drop fuzzy shadow,#1}

% Information box ======================================================
\newtcolorbox{info}[1][]{enhanced,
before skip=5mm,after skip=6mm, % [V3 ADDITION] was 2mm/3mm
boxrule=1pt,left=11mm,right=2mm,top=3mm,bottom=3mm,
colback=cyan!20!white,
colframe=cyan!20!black,
sharp corners,rounded corners=southeast,arc is angular,arc=3mm,
underlay={%
\path[fill=tcbcolback!80!black] ([yshift=3mm]interior.south east)--++(-0.4,-0.1)--++(0.1,-0.2);
\path[draw=tcbcolframe,shorten <=-0.05mm,shorten >=-0.05mm] ([yshift=3mm]interior.south east)--++(-0.4,-0.1)--++(0.1,-0.2);
\path[fill=cyan!50!black,draw=none] (interior.south west) rectangle node[white]{\Huge\bfseries i} ([xshift=9mm]interior.north west);
},
drop fuzzy shadow,#1}

% Mistake box ======================================================
\newtcolorbox{mistake}[1][]{enhanced,
before skip=5mm,after skip=6mm, % [V3 ADDITION] was 2mm/3mm
boxrule=1pt,left=11mm,right=2mm,top=3mm,bottom=3mm,
colback=red!20!white,
colframe=red!20!black,
sharp corners,rounded corners=southeast,arc is angular,arc=3mm,
underlay={%
\path[fill=tcbcolback!80!black] ([yshift=3mm]interior.south east)--++(-0.4,-0.1)--++(0.1,-0.2);
\path[draw=tcbcolframe,shorten <=-0.05mm,shorten >=-0.05mm] ([yshift=3mm]interior.south east)--++(-0.4,-0.1)--++(0.1,-0.2);
\path[fill=red!50!black,draw=none] (interior.south west) rectangle node[white]{\Huge\bfseries $\times$} ([xshift=9mm]interior.north west);
},
drop fuzzy shadow,#1}



% [V3 ADDITION] Worked example box ==================================
% Own counter, numbered within the section (Example 11.2.1, 11.2.2, ...).
% For per-subsection numbering (11.2.1.1) change `section' to `subsection'. Blue `question' boxes are kept for thought-provoking questions
% only; worked examples always use this box. Usage:
%   \begin{example} ...question text... \end{example}
%   \solution{...}
\newtcolorbox[auto counter,number within=section]{example}[1][]{%
enhanced, breakable=false, before skip=5mm, after skip=4mm,
colback=teal!5!white, colframe=teal!55!black, boxrule=0.8pt,
left=3mm, right=3mm, top=2mm, bottom=2mm,
sharp corners, rounded corners=northwest, arc=3mm,
fonttitle=\bfseries, coltitle=white,
title=Example~\thetcbcounter, attach boxed title to top left={yshift=-2mm, xshift=3mm},
boxed title style={colback=teal!55!black, sharp corners, rounded corners=northwest, arc=2mm},
drop fuzzy shadow, #1}

% [V3 ADDITION] Spacing between sections, subsections, parts, paragraphs =
% (fncychap's chapter heading is untouched; only \section and below.)
\makeatletter
\renewcommand\section{\@startsection{section}{1}{\z@}%
  {-2ex \@plus -.5ex \@minus -.2ex}{3.2ex \@plus .3ex}%
  {\normalfont\LARGE\bfseries\color{blue!50!black}}}
\renewcommand\subsection{\@startsection{subsection}{2}{\z@}%
  {-1.5ex \@plus -.5ex \@minus -.2ex}{2.6ex \@plus .3ex}%
  {\normalfont\Large\bfseries}}
\renewcommand\subsubsection{\@startsection{subsubsection}{3}{\z@}%
  {-4ex \@plus -1ex \@minus -.5ex}{1.6ex \@plus .2ex}%
  {\normalfont\large\bfseries\color{teal!50!black}}}
\makeatother
\setlength{\parindent}{0pt}
\setlength{\parskip}{0.9ex plus 0.3ex minus 0.2ex}

% Headers and Footers
\usepackage{fancyhdr,lastpage}
% [V3 R2] grey left-margin line on the out-of-syllabus pages (item 13)
\usepackage{eso-pic}
\newcommand{\oosmargin}{\AddToShipoutPictureBG{\put(12mm,0){\color{gray!75}\rule{1.6pt}{\paperheight}}}}
\pagestyle{fancy}
\lhead{Maximus Pung Jun}
% [BUILD COPY CHANGE 2] Header bug fix: \chaptername is the word "Chapter",
% so the original header printed "Chapter 1: Chapter". \chaptermark now stores
% the bare chapter title in \leftmark, and the header uses \leftmark.
\renewcommand{\chaptermark}[1]{\markboth{#1}{}}
\rhead{Chapter \thechapter: \leftmark}
% \cfoot{Page \thepage\  of \pageref{LastPage}} % buggy, does not appear on chapter pages.

% [V3 R2] tick and cross markers for mistake boxes (G3)
\newcommand{\wrong}{\textcolor{red!70!black}{\ensuremath{\boldsymbol{\times}}}}
\newcommand{\cmark}{\textcolor{green!50!black}{\ensuremath{\checkmark}}}

\newcommand{\xmark}{\wrong}   % alias: red cross
% [V3 R2] out-of-syllabus helpers (DE notes item 13)
%   \oosmargin    : grey line in the left margin on every following page
%   \oosmarginoff : stop the margin line
%   \oosnote      : grey banner "Out of syllabus. For your interest only."
%   outofsyllabus : environment = new page + margin line + banner; ends with
%                   \clearpage and switches the margin line off again
\newcommand{\oosmarginoff}{\ClearShipoutPictureBG}
\newcommand{\oosnote}[1][Nothing in this section is examined in H2 Mathematics.]{\par\noindent\colorbox{gray!15}{\parbox{\dimexpr\textwidth-2\fboxsep}{\textbf{Out of syllabus. For your interest only.} #1}}\par}
\newenvironment{outofsyllabus}{\clearpage\oosmargin}{\clearpage\oosmarginoff}

% Notation as in Maximus's notes: italic d, display style everywhere (G1).
% Fractions inside exponents stay small, so never use these macros inside ^{}.
\newcommand{\dydx}{\dfrac{dy}{dx}}
\newcommand{\dvdx}{\dfrac{dv}{dx}}
\newcommand{\dudx}{\dfrac{du}{dx}}
\newcommand{\ddydx}{\dfrac{d^{2}y}{dx^{2}}}
\newcommand{\ddt}[1]{\dfrac{d#1}{dt}}
% Small syllabus flag (no longer used in the DE notes, kept for other chapters)
\newcommand{\syllabusflag}[1]{\par\noindent{\footnotesize\color{red!70!black}\textbf{Syllabus note.} #1}\par}


% Teacher's copy toggle ================================================
% [BUILD COPY CHANGE 3] \ifx compares the whole word (the original
% \if\conform test was always false). DE_notes_student.tex predefines
% \conform as HideSolutions.
% A student copy defines \conform as HideSolutions BEFORE loading this preamble.
\providecommand\conform{ShowSolutions}
\newcommand\ShowSolutionsKey{ShowSolutions}
% [V3 ADDITION] In the student copy the hidden solution leaves stretchable
% blank space (\vfill), so the free space on each page is shared out as
% working space below the examples instead of collecting at the bottom.
\newcommand\solution[1]{\ifx\conform\ShowSolutionsKey
  \par\noindent\textbf{Solution.} #1\par\medskip\else\par\vspace{\workspace}\fi}
% [V3 R2] minimum working space left under each example in the student copy
\newlength\workspace\setlength\workspace{3cm plus 1fill}


\newcommand{\NN}{\mathbb{N}}
\newcommand{\ZZ}{\mathbb{Z}}
\newcommand{\QQ}{\mathbb{Q}}
\newcommand{\RR}{\mathbb{R}}
\newcommand{\CC}{\mathbb{C}}


% =====================================================================
\usepackage{pifont}
% preamble_local.tex: extra packages and styles for the Vectors chapter
% ONLY. Loaded after preamble_build.tex (a verbatim copy of the shared
% book preamble from the PnC / DE build, so that it can be re-synced by
% copying it over again).
% =====================================================================
\usepackage{array, booktabs}
\usetikzlibrary{arrows.meta, calc, positioning, 3d, decorations.pathreplacing, angles, quotes}
% =====================================================================
% vec_common.tex: shared by the notes (via notes/preamble_local.tex) and
% the slides (via slides/slides_local.tex) for the Vectors chapter.
%
%   CHAPTER NUMBER: change the 13 below and rebuild. The notes set the
%   chapter counter from it, and every slide frame title, Example/Quiz
%   number and box number on the slides uses \chapnum.
% =====================================================================
\def\chapnum{13}

% ---- answers -------------------------------------------------------------
% \ans{key}{value} just prints the value. checks/verify.py reads every
% \ans in the notes and slides and checks it with sympy.
\newcommand{\ans}[2]{#2}

% Column vector, the form used in MF27 and in the examples.
% pmatrix cells are text style, so \dfrac alone still looks small and the
% rows collide. Force display style in every cell and open the rows up.
% Each cell gets a strut as tall as its own content plus \colpad above and
% below, so rows of \dfrac never touch, while integer rows stay compact.
\newlength{\colpad}\setlength{\colpad}{0.55ex}
\newcommand{\colcell}[1]{%
  \setbox0=\hbox{$\displaystyle #1$}%
  \rule[-\dimexpr\dp0+\colpad\relax]{0pt}{\dimexpr\ht0+\dp0+2\colpad\relax}%
  \box0}
\newcommand{\col}[3]{%
  \begingroup
  \renewcommand{\arraystretch}{1.35}%
  \begin{pmatrix}
    \colcell{#1}\\
    \colcell{#2}\\
    \colcell{#3}
  \end{pmatrix}%
  \endgroup
}

% \place{Name}{x}{y}{z} defines a TikZ coordinate. The parentheses are
% inside the macro, so TikZ sees them only after expansion, and
% checks/verify.py can read the drawn coordinates back out.
\newcommand{\place}[4]{\coordinate (#1) at (#2,#3,#4);}

% ---- 3D pictures ----------------------------------------------------------
% Oblique axes: x to the right and slightly down, y out of the page,
% z straight up. Use scale to fit the frame.
\tikzset{
  threeD/.style={
    x={(0.78cm,-0.20cm)},
    y={(0.36cm,0.26cm)},
    z={(0cm,0.78cm)},
    line cap=round, line join=round,
    >=Stealth
  },
  varr/.style={->, thick},
  pt/.style={circle, fill=black, inner sep=1.15pt},
  plane/.style={fill=blue!12, draw=blue!45!black, thick, fill opacity=0.85},
  hid/.style={dashed, black!55}
}
% vecfig.tex: shared 3D diagram style (tikz-3dplot) for Vectors I / II notes and slides.
\usepackage{tikz-3dplot}
\tdplotsetmaincoords{62}{55}
\definecolor{vA}{RGB}{31,90,180}   % first vector / line 1
\definecolor{vB}{RGB}{210,90,20}   % second vector / line 2
\definecolor{vC}{RGB}{0,128,110}   % result / third
\definecolor{vN}{RGB}{140,40,150}  % normals
\definecolor{vP}{RGB}{200,215,240} % plane 1
\definecolor{vQ}{RGB}{245,215,190} % plane 2
\definecolor{vR}{RGB}{200,235,215} % plane 3
\tikzset{
  v3/.style={tdplot_main_coords, line cap=round, line join=round, >=Stealth, scale=#1},
  v3/.default=1,
  ax/.style={->, black!50, thin},
  va/.style={->, very thick, vA},
  vb/.style={->, very thick, vB},
  vc/.style={->, very thick, vC},
  vn/.style={->, very thick, vN},
  ln/.style={very thick},
  dsh/.style={dashed, black!55},
  pp/.style={fill=vP, fill opacity=0.55, draw=vA!70!black, thin},
  pq/.style={fill=vQ, fill opacity=0.55, draw=vB!70!black, thin},
  pr/.style={fill=vR, fill opacity=0.55, draw=vC!70!black, thin},
  dot/.style={circle, fill=black, inner sep=1.2pt},
}
% \vaxes{xlen}{ylen}{zlen}
\newcommand{\vaxes}[3]{%
  \draw[ax] (0,0,0) -- (#1,0,0) node[anchor=west]{$x$};
  \draw[ax] (0,0,0) -- (0,#2,0) node[anchor=south west]{$y$};
  \draw[ax] (0,0,0) -- (0,0,#3) node[anchor=south]{$z$};
  \node[anchor=east] at (0,0,0) {\scriptsize $O$};}
% \rightmark{corner}{along1}{along2} small right-angle square
\newcommand{\rtang}[3]{\draw[thin] ($(#1)!0.25cm!(#2)$) -- ($($(#1)!0.25cm!(#2)$)+($(#1)!0.25cm!(#3)$)-(#1)$) -- ($(#1)!0.25cm!(#3)$);}
\newenvironment{vis}{\par\medskip\begin{center}}{\end{center}\medskip}
% \vang{vertex}{p1}{p2}{label}: small angle arc drawn on the page
\newcommand{\vang}[4]{%
  \coordinate (vg1) at ($(#1)!0.55cm!(#2)$); \coordinate (vg2) at ($(#1)!0.55cm!(#3)$);
  \coordinate (vgm) at ($(vg1)!0.5!(vg2)$);
  \draw[vC, thick] (vg1) .. controls ($(#1)!0.72cm!(vgm)$) .. (vg2);
  \node[vC] at ($(#1)!0.95cm!(vgm)$) {#4};}


% A little air between the parts of an example, and between lines of a list.
\setlength{\itemsep}{0.85ex}
\setlength{\parsep}{0.35ex}

% Student copy: minimum working space under each example. 2.2cm (instead of
% the book's 3cm) is meant to keep the student copy on the same page breaks
% as the teacher copy.
\setlength\workspace{2.2cm plus 1fill}

% =====================================================================
% prelim_fonts.tex: fonts of the real 9758 prelim papers (Word + MathType; pdffonts on
% ACJC 2026 P1/P2, ACJC 2025 P1, JPJC 2026, RVHS 2026, CJC/DHS/TJC/NYJC/HCI/EJC 2025):
%   text Times New Roman; cover/footer Arial; RI page footer Calibri (here Carlito);
%   maths = MathType: Times New Roman letters/digits (italic variables; upright digits,
%   function names, d, e, f), SymbolMT for the integral, sigma, pi and = + < >,
%   scripts 7 pt on 12 pt (about 58 %), tight operator spacing.
% Candidate C (teacher-approved 13 Oct 2026): the real Times New Roman and Arial
% (ttf-mscorefonts-installer) + URW Standard Symbols PS (fonts-urw-base35, a SymbolMT
% clone with the same Symbol encoding) for the integral, sigma, pi, = + < >;
% TeX Gyre Termes Math only for what MathType draws itself (radicals, stretchy brackets,
% fraction bars, arrows). House style: upright d (\mathrm{d}x), upright e and f.
% Loaded once by di_style / de_core (HW sheets) and prelim_style (Tests, solutions).
% =====================================================================
\ifdefined\prelimfontsloaded\endinput\fi
\def\prelimfontsloaded{1}
\edef\pfrestoreat{\catcode`\noexpand\@=\the\catcode`\@\relax}\makeatletter
% MathType script sizes: 7 pt superscripts on 12 pt text (58 %), 5 pt second level
\DeclareMathSizes{12}{12}{7}{5}
\DeclareMathSizes{11}{11}{6.5}{5}
\DeclareMathSizes{10.95}{10.95}{6.5}{5}
\DeclareMathSizes{10}{10}{6}{5}
\DeclareMathSizes{9}{9}{5.5}{5}
\DeclareMathSizes{14.4}{14.4}{8.5}{6}
\setmainfont{Times New Roman}
\setsansfont{Arial}
\newfontfamily\coverfont{Arial}
\newfontfamily\calibrifont{Carlito}
\usepackage{unicode-math}
\setmathfont{TeX Gyre Termes Math}
\setmathfont[range={up/{latin,Latin,num}}]{Times New Roman}
\setmathfont[range={it/{latin,Latin}}]{Times New Roman Italic}
\setmathfont[range={bfup/{latin,Latin,num}}]{Times New Roman Bold}
\setmathfont[range={"002B,"003D,"003C,"003E}]{Standard Symbols PS}
\newfontfamily\pf@sym{Standard Symbols PS}
\newfontfamily\pf@symi[FakeSlant=0.21]{Standard Symbols PS}
% MathType-like spacing (tighter than TeX around = + -)
\AtBeginDocument{\thickmuskip=2.5mu plus 1mu\medmuskip=1.5mu plus 1mu\thinmuskip=3mu
% Symbol integral and sigma (18 pt at 12 pt text, as MathType), slanted Symbol pi
 \renewcommand{\int}{\mathop{\vcenter{\hbox{\pf@sym\fontsize{\fpeval{1.5*\f@size}}{0pt}\selectfont\char"F2}}}\nolimits}%
 \renewcommand{\sum}{\mathop{\vcenter{\hbox{\pf@sym\fontsize{\fpeval{1.5*\f@size}}{0pt}\selectfont\char"E5}}}\limits}%
 \renewcommand{\pi}{\mathord{\text{\pf@symi\char"70}}}}
\pfrestoreat

\usepackage{needspace,longtable,multicol,enumitem,array}
\providecommand{\vv}[1]{\overrightarrow{#1}}\providecommand{\cvv}[3]{\begin{pmatrix}#1\\#2\\#3\end{pmatrix}}
\raggedbottom
\setlength{\parindent}{0pt}

\fancyhf{}\lhead{\small Vectors I: HW3 Solutions}\rhead{\small Compiled by Maximus Pung Jun}
\cfoot{\small\thepage}
\renewcommand{\headrulewidth}{0.4pt}
\newcounter{wq}
\newcommand{\Q}[1]{\par\needspace{5\baselineskip}\medskip\stepcounter{wq}\noindent\textbf{\thewq}\quad #1\par}
\begin{document}
\thispagestyle{fancy}
{\LARGE\scshape Vectors I Homework 3: Vector Product --- Solutions}\par\smallskip
{Compiled by Maximus Pung Jun}\par\smallskip
{\small\textcolor{black!60}{Official worked solutions for sourced past-paper questions; True/False answers from the Basics Worksheet.}}\par\medskip
\Q{For the following statements, state whether it is \textbf{true} or \textbf{false}.}
\par\smallskip{\setlength{\extrarowheight}{6pt}\renewcommand{\arraystretch}{1.18}\begin{longtable}{|>{\centering\arraybackslash}p{0.9cm}|p{\dimexpr\linewidth-0.9cm-1.6cm-6\tabcolsep-4\arrayrulewidth\relax}|>{\centering\arraybackslash}p{1.6cm}|}
\hline &\textbf{Statement}&\textbf{T/F}\\\hline\endhead
(a)&$\mathbf a\times\mathbf a=\mathbf 0$.\par\smallskip $\displaystyle |\mathbf a\times\mathbf a|=|\mathbf a|^2\sin0^\circ=0$\rule[-6pt]{0pt}{6pt}&\textbf{T}\\\hline
(b)&$\mathbf a\times\mathbf a=0$.\par\smallskip $\displaystyle \mathbf a\times\mathbf a=\mathbf 0\neq0\ \text{(left side is the zero vector)}$\rule[-6pt]{0pt}{6pt}&\textbf{F}\\\hline
(c)&$\mathbf i\times\mathbf j=\mathbf k$.\par\smallskip $\displaystyle \cvv{1}{0}{0}\times\cvv{0}{1}{0}=\cvv{0}{0}{1}$\rule[-6pt]{0pt}{6pt}&\textbf{T}\\\hline
(d)&$\mathbf a\times\mathbf b=\mathbf b\times\mathbf a$.\par\smallskip $\displaystyle \mathbf i\times\mathbf j=\mathbf k,\ \mathbf j\times\mathbf i=-\mathbf k$\rule[-6pt]{0pt}{6pt}&\textbf{F}\\\hline
(e)&$\mathbf a\times\mathbf b=-(\mathbf b\times\mathbf a)$.\par\smallskip $\displaystyle \mathbf a\times\mathbf b=-\mathbf b\times\mathbf a$\rule[-6pt]{0pt}{6pt}&\textbf{T}\\\hline
(f)&$\mathbf a\times\mathbf b$ is a scalar.\par\smallskip $\displaystyle \mathbf i\times\mathbf j=\mathbf k\ \text{(a vector)}$\rule[-6pt]{0pt}{6pt}&\textbf{F}\\\hline
(g)&$\mathbf k\times\mathbf j=\mathbf i$.\par\smallskip $\displaystyle \mathbf j\times\mathbf k=\mathbf i\ \therefore\ \mathbf k\times\mathbf j=-\mathbf i$\rule[-6pt]{0pt}{6pt}&\textbf{F}\\\hline
(h)&$\cvv{1}{2}{3}\times\cvv{2}{4}{6}=\mathbf 0$.\par\smallskip $\displaystyle \cvv{2}{4}{6}=2\cvv{1}{2}{3}$\rule[-6pt]{0pt}{6pt}&\textbf{T}\\\hline
(i)&$\mathbf a\times(\mathbf b+\mathbf c)=\mathbf a\times\mathbf b+\mathbf a\times\mathbf c$.\par\smallskip $\displaystyle \text{Distributive law}$\rule[-6pt]{0pt}{6pt}&\textbf{T}\\\hline
(j)&$(\lambda\mathbf a)\times\mathbf b=\lambda(\mathbf a\times\mathbf b)$ for all $\lambda\in\mathbb R$.\par\smallskip $\displaystyle (\lambda\mathbf a)\times\mathbf b=\mathbf a\times(\lambda\mathbf b)=\lambda(\mathbf a\times\mathbf b)$\rule[-6pt]{0pt}{6pt}&\textbf{T}\\\hline
(k)&$\mathbf a\times\mathbf b$ is perpendicular to both $\mathbf a$ and $\mathbf b$.\par\smallskip $\displaystyle (\mathbf a\times\mathbf b)\cdot\mathbf a=0,\ (\mathbf a\times\mathbf b)\cdot\mathbf b=0$\rule[-6pt]{0pt}{6pt}&\textbf{T}\\\hline
(l)&$(\mathbf a\times\mathbf b)\cdot\mathbf c$ is a scalar.\par\smallskip $\displaystyle (\mathbf a\times\mathbf b)\cdot\mathbf c=\text{vector}\cdot\text{vector}$\rule[-6pt]{0pt}{6pt}&\textbf{T}\\\hline
(m)&$|\mathbf a\times\mathbf b|=|\mathbf a||\mathbf b|\sin\theta$, where $\theta$ is the angle between $\mathbf a$ and $\mathbf b$.\par\smallskip $\displaystyle \text{Definition of vector product}$\rule[-6pt]{0pt}{6pt}&\textbf{T}\\\hline
(n)&$|\mathbf a\times\mathbf b|=|\mathbf a||\mathbf b|$.\par\smallskip $\displaystyle |\mathbf i\times\mathbf i|=0\neq1$\rule[-6pt]{0pt}{6pt}&\textbf{F}\\\hline
(o)&If $\mathbf a\times\mathbf b=\mathbf 0$, then $\mathbf a=\mathbf 0$ or $\mathbf b=\mathbf 0$.\par\smallskip $\displaystyle \mathbf i\times2\mathbf i=\mathbf 0,\ \mathbf i\neq\mathbf 0,\ 2\mathbf i\neq\mathbf 0$\rule[-6pt]{0pt}{6pt}&\textbf{F}\\\hline
(p)&For non-zero $\mathbf a$ and $\mathbf b$, $\mathbf a\times\mathbf b=\mathbf 0$ means that $\mathbf a$ is parallel to $\mathbf b$.\par\smallskip $\displaystyle |\mathbf a||\mathbf b|\sin\theta=0\ \therefore\ \theta=0^\circ\ \text{or}\ 180^\circ$\rule[-6pt]{0pt}{6pt}&\textbf{T}\\\hline
(q)&The area of the parallelogram with adjacent sides $\mathbf a$ and $\mathbf b$ is $|\mathbf a\times\mathbf b|$.\par\smallskip $\displaystyle \text{Area}=|\mathbf a||\mathbf b|\sin\theta=|\mathbf a\times\mathbf b|$\rule[-6pt]{0pt}{6pt}&\textbf{T}\\\hline
(r)&The area of triangle $OAB$ is $|\vv{OA}\times\vv{OB}|$.\par\smallskip $\displaystyle \text{Area}=\tfrac12|\vv{OA}\times\vv{OB}|$\rule[-6pt]{0pt}{6pt}&\textbf{F}\\\hline
(s)&$\mathbf a\times(\mathbf b\cdot\mathbf c)$ is a vector.\par\smallskip $\displaystyle \mathbf b\cdot\mathbf c\ \text{is a scalar}\ \therefore\ \text{meaningless}$\rule[-6pt]{0pt}{6pt}&\textbf{F}\\\hline
(t)&For non-zero $\mathbf a$ and $\mathbf b$, if $|\mathbf a\times\mathbf b|=|\mathbf a||\mathbf b|$, then $\mathbf a$ is perpendicular to $\mathbf b$.\par\smallskip $\displaystyle \sin\theta=1\ \therefore\ \theta=90^\circ$\rule[-6pt]{0pt}{6pt}&\textbf{T}\\\hline
(u)&$(\mathbf a\times\mathbf b)\cdot\mathbf c=\mathbf a\cdot(\mathbf b\times\mathbf c)$.\par\smallskip $\displaystyle \text{both}=\det\left(\mathbf a\ \mathbf b\ \mathbf c\right)$\rule[-6pt]{0pt}{6pt}&\textbf{T}\\\hline
(v)&$|\mathbf a\times\mathbf b|^2+(\mathbf a\cdot\mathbf b)^2=|\mathbf a|^2|\mathbf b|^2$.\par\smallskip $\displaystyle |\mathbf a|^2|\mathbf b|^2(\sin^2\theta+\cos^2\theta)=|\mathbf a|^2|\mathbf b|^2$\rule[-6pt]{0pt}{6pt}&\textbf{T}\\\hline
(w)&$\mathbf a\times\mathbf b=\mathbf a\times\mathbf c$ means that $\mathbf b=\mathbf c$.\par\smallskip $\displaystyle \mathbf i\times\mathbf j=\mathbf k=\mathbf i\times(\mathbf i+\mathbf j),\ \mathbf j\neq\mathbf i+\mathbf j$\rule[-6pt]{0pt}{6pt}&\textbf{F}\\\hline
(x)&$\mathbf a\times(\mathbf b\times\mathbf c)=(\mathbf a\times\mathbf b)\times\mathbf c$.\par\smallskip $\displaystyle \mathbf i\times(\mathbf i\times\mathbf j)=\mathbf i\times\mathbf k=-\mathbf j,\ (\mathbf i\times\mathbf i)\times\mathbf j=\mathbf 0$\rule[-6pt]{0pt}{6pt}&\textbf{F}\\\hline
(y)&The shortest distance from $B$ to the line through $A$ parallel to $\mathbf d$ is $|\vv{AB}\times\mathbf d|$.\par\smallskip $\displaystyle \text{distance}=|\vv{AB}\times\hat{\mathbf d}|=\dfrac{|\vv{AB}\times\mathbf d|}{|\mathbf d|}$\rule[-6pt]{0pt}{6pt}&\textbf{F}\\\hline
(z)&If $\mathbf a\times\mathbf b=\mathbf 0$ and $\mathbf a\cdot\mathbf b=0$, then $\mathbf a=\mathbf 0$ or $\mathbf b=\mathbf 0$.\par\smallskip $\displaystyle |\mathbf a|^2|\mathbf b|^2=|\mathbf a\times\mathbf b|^2+(\mathbf a\cdot\mathbf b)^2=0$\rule[-6pt]{0pt}{6pt}&\textbf{T}\\\hline\end{longtable}}
\clearpage
\par\noindent\begin{minipage}{\linewidth}
% Full selected questions, in owner-selected candidate order: 6,5,2,4,9,12,11.
\Q{\textcolor{black!60}{\small[JPJC 2024 Promo Q11]}\par\smallskip
Referred to the origin $O$, points $A$ and $B$ have position vectors $\mathbf a$ and $\mathbf b$ respectively, where $\mathbf a$ and $\mathbf b$ are non-zero and non-parallel vectors. Point $C$ is the midpoint of $OA$ and point $D$ lies on $BC$ such that $BD:DC=2:3$.
\begin{enumerate}[label=\textbf{(\roman*)},leftmargin=2.2em,itemsep=3pt]
\item Find the position vectors $\vv{OC}$ and $\vv{OD}$, giving your answers in terms of $\mathbf a$ and $\mathbf b$.\hfill[2]
\item Find the area of triangle $OCD$ in the form of $k|\mathbf a\times\mathbf b|$, where $k$ is a constant to be found.\hfill[3]
\item Given that $|\mathbf a|=2$, $|\mathbf b|=\sqrt2$ and angle $AOB$ is $\dfrac\pi4$ radian, using scalar product, show that $CD$ is perpendicular to $OA$.\hfill[3]
\item Given instead that $\mathbf a$ is a unit vector, give a geometrical interpretation of $|\mathbf a\cdot\mathbf b|$.\hfill[1]
\end{enumerate}}
\end{minipage}\par
\par\medskip\hrule\medskip
% Faithful equation transcriptions of selected official marking-scheme methods.
\textbf{Solution}\par\smallskip
\textbf{(i)} Given $\vv{OA}=\mathbf a$, $\vv{OB}=\mathbf b$, $\vv{OC}=\dfrac12\mathbf a$.
\begin{align*}
\vv{OD}&=\frac{2\mathbf c+3\mathbf b}{5}
=\frac15\left[2\left(\frac12\mathbf a\right)+3\mathbf b\right]
=\frac15(\mathbf a+3\mathbf b).
\end{align*}
\textbf{(ii)}
\begin{align*}
\text{Area of triangle }OCD
&=\frac12|\vv{OC}\times\vv{OD}|\\
&=\frac12\left|\frac12\mathbf a\times\frac15(\mathbf a+3\mathbf b)\right|\\
&=\frac12\left|\frac1{10}(\mathbf a\times\mathbf a)+\frac3{10}(\mathbf a\times\mathbf b)\right|\\
&=\frac12\left|\frac3{10}(\mathbf a\times\mathbf b)\right|\qquad(\mathbf a\times\mathbf a=\mathbf0)\\
&=\frac3{20}|\mathbf a\times\mathbf b|,\qquad k=\frac3{20}.
\end{align*}
\textbf{(iii)}
\begin{align*}
\vv{CD}&=\vv{OD}-\vv{OC}=\frac15(\mathbf a+3\mathbf b)-\frac12\mathbf a
=\frac35\mathbf b-\frac3{10}\mathbf a\\
\vv{CD}\cdot\vv{OA}
&=\left(\frac35\mathbf b-\frac3{10}\mathbf a\right)\cdot\mathbf a\\
&=\frac35\mathbf b\cdot\mathbf a-\frac3{10}\mathbf a\cdot\mathbf a\\
&=\frac35|\mathbf b||\mathbf a|\cos\frac\pi4-\frac3{10}|\mathbf a|^2\\
&=\frac35(\sqrt2)(2)\frac{\sqrt2}{2}-\frac3{10}(2)^2
=\frac65-\frac{12}{10}=0.
\end{align*}
Since $\vv{CD}\cdot\vv{OA}=0$, $CD$ is perpendicular to $OA$. (Shown)\par
\textbf{(iv)} $|\mathbf a\cdot\mathbf b|$ is the length of projection of $\vv{OB}$ onto $\vv{OA}$ (or a line or vector parallel to $\vv{OA}$).
\clearpage
\par\noindent\begin{minipage}{\linewidth}
\Q{\textcolor{black!60}{\small[MI 2024 Promo Q4]}\par\smallskip
Referred to the origin $O$, the points $A$, $B$ and $C$ have position vectors $\mathbf a$, $\mathbf b$ and $\mathbf c$. $C$ lies on $AB$ such that $AC:AB=1:3$.
\begin{enumerate}[label=\textbf{(\roman*)},leftmargin=2.2em,itemsep=3pt]
\item Find $\mathbf c$ in terms of $\mathbf a$ and $\mathbf b$.\hfill[2]
\item Show that the area of triangle $OAC$ is given by $k|\mathbf a\times\mathbf b|$, where $k$ is a constant to be determined.\hfill[2]
\end{enumerate}
It is now given that $\mathbf a$ and $\mathbf b$ are unit vectors.
\begin{enumerate}[label=\textbf{(\roman*)},leftmargin=2.2em,itemsep=3pt,start=3]
\item Using the result in part (i), find the value of $(2\mathbf a-\mathbf b)\cdot\mathbf c$.\hfill[3]
\end{enumerate}}
\end{minipage}\par
\par\medskip\hrule\medskip
\textbf{Solution}\par\smallskip
\textbf{(i)}
\begin{align*}
AC:AB&=1:3\quad\Rightarrow\quad AC:CB=1:2\\
\mathbf c&=\frac{(1)\mathbf b+(2)\mathbf a}{1+2}=\frac{2\mathbf a+\mathbf b}{3}.
\end{align*}
\textbf{(ii)}
\begin{align*}
\text{Area of triangle }OAC
&=\frac13\text{ Area of triangle }OAB\\
&=\frac13\times\frac12|\vv{OA}\times\vv{OB}|\\
&=\frac16|\mathbf a\times\mathbf b|,\qquad k=\frac16.
\end{align*}
\textbf{(iii)}
\begin{align*}
(2\mathbf a-\mathbf b)\cdot\mathbf c
&=(2\mathbf a-\mathbf b)\cdot\left(\frac{2\mathbf a+\mathbf b}{3}\right)\\
&=\frac13(2\mathbf a-\mathbf b)\cdot(2\mathbf a+\mathbf b)\\
&=\frac13(4\mathbf a\cdot\mathbf a+2\mathbf a\cdot\mathbf b-2\mathbf b\cdot\mathbf a-\mathbf b\cdot\mathbf b)\\
&=\frac13(4|\mathbf a|^2-|\mathbf b|^2)\\
&=\frac13\left(4(1)^2-(1)^2\right)\quad\text{since $\mathbf a$, $\mathbf b$ are unit vectors}\\
&=1.
\end{align*}
\clearpage
\par\noindent\begin{minipage}{\linewidth}
\Q{\textcolor{black!60}{\small[ASRJC 2024 Promo Q5]}\par\smallskip
In triangle $PQR$, the vectors $\mathbf a$, $\mathbf b$ and $\mathbf c$ represent the vectors $\vv{PQ}$, $\vv{QR}$ and $\vv{PR}$ respectively.
\begin{center}\begin{tikzpicture}[scale=0.85,font=\small,>=Stealth]
\coordinate(P)at(0,0);\coordinate(Q)at(5,0);\coordinate(R)at(4.15,2.2);
\draw[semithick](P)--(Q)--(R)--cycle;
\draw[->,semithick](P)--(2.7,0);\draw[->,semithick](Q)--(4.57,1.12);\draw[->,semithick](P)--(2.3,1.219);
\node[below left]at(P){$P$};\node[below right]at(Q){$Q$};\node[above]at(R){$R$};
\node[below]at(2.6,0){$\mathbf a$};\node[right]at(4.7,1.15){$\mathbf b$};\node[above left]at(2.1,1.1){$\mathbf c$};
\end{tikzpicture}\end{center}
\begin{enumerate}[label=\textbf{(\alph*)},leftmargin=2.2em,itemsep=3pt]
\item State a geometrical interpretation of $|\mathbf c\times\hat{\mathbf a}|$ and hence show that the area of the triangle $PQR$ can be written as $\dfrac12|\mathbf c\times\mathbf a|$ units$^2$.\hfill[3]
\item With clear explanations, show that $|\mathbf c\times\mathbf a|=|\mathbf b\times\mathbf a|$ and deduce that
\[\frac{|\mathbf c|}{\sin\angle PQR}=\frac{|\mathbf b|}{\sin\angle QPR}\quad\text{(Sine Rule).}\]\hfill[2]
\end{enumerate}}
\end{minipage}\par
\par\medskip\hrule\medskip
\textbf{Solution}\par\smallskip
\textbf{(a)} $|\mathbf c\times\hat{\mathbf a}|$ represents the perpendicular distance from the point $R$ to the line $PQ$.
\begin{align*}
\text{Area of }\triangle PQR&=\frac12\times\text{base}\times\text{height}\\
&=\frac12|\mathbf a|\,|\mathbf c\times\hat{\mathbf a}|\\
&=\frac12|\mathbf a|\left|\mathbf c\times\frac{\mathbf a}{|\mathbf a|}\right|\\
&=\frac12|\mathbf c\times\mathbf a|\text{ units}^2.\quad\text{(Shown)}
\end{align*}
\textbf{(b)} By replacing $\mathbf a$ with $-\mathbf a$ and $\mathbf c$ with $\mathbf b$ in part (a),
\begin{align*}
\text{Area of }\triangle PQR&=\frac12|\mathbf b\times(-\mathbf a)|=\frac12|\mathbf b\times\mathbf a|\\
\frac12|\mathbf c\times\mathbf a|&=\frac12|\mathbf b\times\mathbf a|\\
|\mathbf c\times\mathbf a|&=|\mathbf b\times\mathbf a|\quad\text{(Shown)}\\
|\mathbf c||\mathbf a|\sin\angle QPR&=|\mathbf b||\mathbf a|\sin\angle PQR\\
\frac{|\mathbf c|}{\sin\angle PQR}&=\frac{|\mathbf b|}{\sin\angle QPR}.\quad\text{(Shown)}
\end{align*}
\clearpage
\par\noindent\begin{minipage}{\linewidth}
\Q{\textcolor{black!60}{\small[CJC 2024 Promo Q7]}\par\smallskip
\begin{enumerate}[label=\textbf{(\alph*)},leftmargin=2.2em,itemsep=5pt]
\item Three non-zero vectors $\mathbf p$, $\mathbf q$ and $\mathbf r$ are such that $3\mathbf p\times\mathbf q=\mathbf r\times\mathbf p$.\par
Find a linear relationship between $\mathbf p$, $\mathbf q$ and $\mathbf r$.\hfill[3]
\item Referred to the origin $O$, the points $A$ and $B$ have position vectors $\mathbf a$ and $\mathbf b$.\par
It is given that $\mathbf a$ and $\mathbf b$ are non-parallel unit vectors.\par
The point $B$ divides $AC$ in the ratio $1:3$.\par
The point $D$ lies on $OB$ produced such that $OB:OD=1:m$, where $m\in\mathbb Z$, $m\ge2$.\par
Given that $\vv{AB}$ is perpendicular to $\vv{CD}$, find the numerical value of $m$.\hfill[4]
\end{enumerate}}
\end{minipage}\par
\par\medskip\hrule\medskip
\textbf{Solution}\par\smallskip
\textbf{(a)}
\begin{align*}
3\mathbf p\times\mathbf q&=\mathbf r\times\mathbf p\\
(3\mathbf p\times\mathbf q)-(\mathbf r\times\mathbf p)&=\mathbf0\\
(3\mathbf p\times\mathbf q)+(\mathbf p\times\mathbf r)&=\mathbf0\\
\mathbf p\times(3\mathbf q+\mathbf r)&=\mathbf0\\
\mathbf p&\parallel3\mathbf q+\mathbf r\\
\mathbf p&=k(3\mathbf q+\mathbf r),\quad k\in\mathbb R\setminus\{0\}.
\end{align*}
{\small\textit{Verification note (not part of the official scheme):} the printed last line assumes $3\mathbf q+\mathbf r\ne\mathbf0$. The question also permits $3\mathbf q+\mathbf r=\mathbf0$. The unrestricted relationship is $3\mathbf q+\mathbf r=\lambda\mathbf p$, $\lambda\in\mathbb R$, including $\lambda=0$.}\par\smallskip
\textbf{(b)}
\begin{align*}
\vv{OB}&=\frac{3\vv{OA}+\vv{OC}}4\\
\vv{OC}&=4\mathbf b-3\mathbf a\\
\vv{OD}&=m\vv{OB}=m\mathbf b\\
\vv{AB}&=\mathbf b-\mathbf a\\
\vv{CD}&=m\mathbf b-(4\mathbf b-3\mathbf a)=3\mathbf a+(m-4)\mathbf b.
\end{align*}
Given that $\vv{AB}$ is perpendicular to $\vv{CD}$,
\begin{align*}
\vv{AB}\cdot\vv{CD}&=0\\
(\mathbf b-\mathbf a)\cdot[3\mathbf a+(m-4)\mathbf b]&=0\\
3\mathbf a\cdot\mathbf b+(m-4)\mathbf b\cdot\mathbf b-3\mathbf a\cdot\mathbf a-(m-4)\mathbf a\cdot\mathbf b&=0\\
3\mathbf a\cdot\mathbf b+(m-4)|\mathbf b|^2-3|\mathbf a|^2-(m-4)\mathbf a\cdot\mathbf b&=0\\
(m-4)-3-(m-7)\mathbf a\cdot\mathbf b&=0\\
m(1-\mathbf a\cdot\mathbf b)&=3-7\mathbf a\cdot\mathbf b+4\\
m&=\frac{3-7\mathbf a\cdot\mathbf b+4}{1-\mathbf a\cdot\mathbf b}
=\frac{7(1-\mathbf a\cdot\mathbf b)}{1-\mathbf a\cdot\mathbf b}=7.
\end{align*}
\clearpage
\par\noindent\begin{minipage}{\linewidth}
\Q{\textcolor{black!60}{\small[SAJC 2024 Promo Q9]}\par\smallskip
\begin{enumerate}[label=\textbf{(\alph*)},leftmargin=2.2em,itemsep=5pt]
\item Relative to the origin $O$, the points $A$ and $B$ have position vectors $\mathbf a$ and $\mathbf b$ respectively, where $\mathbf a$ and $\mathbf b$ are non-parallel. It is given that $|\mathbf a|=9$, $|\mathbf b|=1$ and $|2\mathbf a+\mathbf b|=2\sqrt{74}$.
\begin{enumerate}[label=(\roman*),leftmargin=1.8em,itemsep=3pt]
\item By considering $(2\mathbf a+\mathbf b)\cdot(2\mathbf a+\mathbf b)$, show that $\mathbf a\cdot\mathbf b=-\dfrac{29}{4}$.\hfill[3]
\item Give a geometrical meaning of $|\mathbf a\cdot\mathbf b|$.\hfill[1]
\item The points $P$, $Q$ and $R$ have position vectors $7\mathbf a-5\mathbf b$, $6\mathbf a+5\mathbf b$ and $9\mathbf a+\lambda\mathbf b$ respectively. Given that these three points are collinear, find the value of $\lambda$.\hfill[3]
\end{enumerate}
\item In the triangle $OUV$ where $O$ is the origin, the position vectors of the points $U$ and $V$ are $\mathbf u$ and $\mathbf v$ respectively. The point $W$ is the midpoint of $OU$, and the point $X$ has position vector given by $\dfrac1{16}\mathbf u+\dfrac34\mathbf v$. Show that the area of triangle $UWX$ can be written as $m|\mathbf u\times\mathbf v|$ where $m$ is a constant to be found.\hfill[3]
\end{enumerate}}
\end{minipage}\par
\par\medskip\hrule\medskip
\textbf{Solution}\par\smallskip
\textbf{(a)(i)}
\begin{align*}
(2\mathbf a+\mathbf b)\cdot(2\mathbf a+\mathbf b)
&=4(\mathbf a\cdot\mathbf a)+2(\mathbf b\cdot\mathbf a)+2(\mathbf a\cdot\mathbf b)+(\mathbf b\cdot\mathbf b)\\
&=4|\mathbf a|^2+4(\mathbf a\cdot\mathbf b)+|\mathbf b|^2\\
|2\mathbf a+\mathbf b|^2&=4|\mathbf a|^2+4(\mathbf a\cdot\mathbf b)+|\mathbf b|^2\\
(2\sqrt{74})^2&=4(9^2)+4(\mathbf a\cdot\mathbf b)+(1^2)\\
296&=4(9^2)+4(\mathbf a\cdot\mathbf b)+(1^2)\\
4(\mathbf a\cdot\mathbf b)&=296-4(9^2)-(1^2)\\
\mathbf a\cdot\mathbf b&=-\frac{29}4.
\end{align*}
\textbf{(a)(ii)} $|\mathbf a\cdot\mathbf b|$ is the length of projection of $\vv{OA}$ onto $\vv{OB}$.\par
\textbf{(a)(iii)} If $P$, $Q$ and $R$ are collinear, then $k\vv{PQ}=\vv{PR}$.
\begin{align*}
k[(6\mathbf a+5\mathbf b)-(7\mathbf a-5\mathbf b)]&=[(9\mathbf a+\lambda\mathbf b)-(7\mathbf a-5\mathbf b)]\\
k[-\mathbf a+10\mathbf b]&=[2\mathbf a+(\lambda+5)\mathbf b].
\end{align*}
Since $\mathbf a$ and $\mathbf b$ are non-parallel, non-zero vectors,
\begin{align*}
k&=-2\\
\lambda+5&=10k=-20\\
\lambda&=-25.
\end{align*}
\textbf{(b)}
\begin{align*}
\vv{UX}&=\vv{OX}-\vv{OU}=\frac1{16}\mathbf u+\frac34\mathbf v-\mathbf u=\frac34\mathbf v-\frac{15}{16}\mathbf u\\
\vv{UW}&=\vv{OW}-\vv{OU}=\frac12\mathbf u-\mathbf u=-\frac12\mathbf u\\
\text{Area of triangle }UWX
&=\frac12|\vv{UW}\times\vv{UX}|\\
&=\frac12\left|\frac12\mathbf u\times\left(\frac34\mathbf v-\frac{15}{16}\mathbf u\right)\right|\\
&=\frac3{16}|\mathbf u\times\mathbf v|\qquad(\mathbf u\times\mathbf u=\mathbf0)\\
m&=\frac3{16}.
\end{align*}
\clearpage
\par\noindent\begin{minipage}{\linewidth}
\Q{\textcolor{black!60}{\small[CJC 2025 Promo Q6]}\par\smallskip
The diagram shows a unit circle with centre $O$ and diameter $AB$. The point $C$ lies on the circumference of the circle. The position vectors of $A$, $B$ and $C$ are $\mathbf a$, $\mathbf b$ and $\mathbf c$ respectively.
\begin{center}\begin{tikzpicture}[scale=.72,font=\small,>=Stealth]
\coordinate(O)at(0,0);\coordinate(A)at(2,0);\coordinate(B)at(-2,0);\coordinate(C)at(-.9,1.786);
\draw[semithick](O)circle(2);\draw[semithick](B)--(A)--(C)--cycle;
\draw[->,semithick](O)--(1.2,0);\draw[->,semithick](O)--(-1.2,0);\draw[->,semithick](O)--(C);
\node[below]at(O){$O$};\node[right]at(A){$A$};\node[left]at(B){$B$};\node[above]at(C){$C$};
\node[below]at(1,0){$\mathbf a$};\node[below]at(-1,0){$\mathbf b$};\node[left]at(-.45,.9){$\mathbf c$};
\end{tikzpicture}\end{center}
\begin{enumerate}[label=\textbf{(\alph*)},leftmargin=2.2em,itemsep=3pt]
\item By considering an appropriate scalar product, show that $AC$ and $BC$ are perpendicular.\hfill[4]
\item Give a geometrical interpretation of $|\mathbf a\cdot(\mathbf c-\mathbf a)|$.\hfill[1]
\end{enumerate}
The point $D$ is such that $ABCD$ is a parallelogram. The point $M$ is the mid-point of $AD$ and the point $R$ lies on $OA$ produced such that $OA:AR=1:2$.
\begin{center}\begin{tikzpicture}[scale=.64,font=\small]
\coordinate(O)at(0,0);\coordinate(A)at(2,0);\coordinate(B)at(-2,0);\coordinate(C)at(-.9,1.786);\coordinate(D)at(3.1,1.786);\coordinate(M)at(2.55,.893);\coordinate(R)at(6,0);
\draw[semithick](O)circle(2);\draw[semithick](B)--(R);\draw[semithick](B)--(C)--(D)--(A)--(C);
\fill(O)circle(1.6pt);\fill(M)circle(1.6pt);
\node[below]at(O){$O$};\node[below]at(A){$A$};\node[left]at(B){$B$};\node[above]at(C){$C$};\node[above]at(D){$D$};\node[right]at(M){$M$};\node[below]at(R){$R$};
\end{tikzpicture}\end{center}
\begin{enumerate}[label=\textbf{(\alph*)},leftmargin=2.2em,itemsep=3pt,start=3]
\item By considering $\vv{CM}\times\vv{CR}$ or otherwise, show that the points $C$, $M$ and $R$ are collinear.\hfill[3]
\end{enumerate}}
\end{minipage}\par
\par\medskip\hrule\medskip
\textbf{Solution}\par\smallskip
\textbf{(a)}
\begin{align*}
\vv{AC}&=\mathbf c-\mathbf a,\qquad\vv{BC}=\mathbf c-\mathbf b\\
\vv{AC}\cdot\vv{BC}
&=(\mathbf c-\mathbf a)\cdot(\mathbf c-\mathbf b)\\
&=(\mathbf c-\mathbf a)\cdot(\mathbf c-(-\mathbf a))\\
&=(\mathbf c-\mathbf a)\cdot(\mathbf c+\mathbf a)\\
&=\mathbf c\cdot\mathbf c-\mathbf a\cdot\mathbf a\\
&=|\mathbf c|^2-|\mathbf a|^2\\
&=0\qquad(|\mathbf c|=|\mathbf a|=\text{radius of circle}).
\end{align*}
Hence $AC$ and $BC$ are perpendicular.\par
\textbf{(b)} $|\mathbf a\cdot(\mathbf c-\mathbf a)|$ is the length of projection of $\vv{AC}$ onto $\vv{OA}$.\par
\textbf{(c)} By Ratio Theorem,
\begin{align*}
\vv{OM}&=\frac{\vv{OD}+\vv{OA}}2\\
&=\frac12(\mathbf c+2\mathbf a+\mathbf a)=\frac12\mathbf c+\frac32\mathbf a\\
\vv{CM}\times\vv{CR}
&=(\vv{OM}-\vv{OC})\times(\vv{OR}-\vv{OC})\\
&=\left[\frac12\mathbf c+\frac32\mathbf a-\mathbf c\right]\times(3\mathbf a-\mathbf c)\\
&=\left(\frac32\mathbf a-\frac12\mathbf c\right)\times(3\mathbf a-\mathbf c)\\
&=\frac92\mathbf a\times\mathbf a-\frac32\mathbf a\times\mathbf c-\frac32\mathbf c\times\mathbf a+\frac12\mathbf c\times\mathbf c\\
&=\mathbf0-\frac32\mathbf a\times\mathbf c+\frac32\mathbf a\times\mathbf c+\mathbf0\\
&=\mathbf0.
\end{align*}
$\vv{CM}$ and $\vv{CR}$ are parallel. Since $CR$ and $CM$ are parallel and $C$ is a common point, hence $C$, $M$ and $R$ are collinear.
\par\vspace{\baselineskip}\begin{center}\textbf{END OF HOMEWORK}\end{center}\end{document}
